\documentclass[11pt,a4paper]{article}
\usepackage{notebooks}
\usepackage{amsmath}
\usepackage{enumitem}
\setlist[itemize]{leftmargin=1.4em,itemsep=4pt,topsep=4pt}
\setlist[enumerate]{leftmargin=1.6em,itemsep=2pt,topsep=4pt}

\newcommand{\valuetable}[1]{%
{\arrayrulecolor{gray!60}\begin{tabularx}{\textwidth}{@{}>{\raggedright\arraybackslash}p{0.30\textwidth}X@{}}
\toprule
\rowcolor{tablehead} \textbf{Value} & \textbf{Meaning} \\
\midrule
#1
\bottomrule
\end{tabularx}}
}

\newcommand{\pyheading}[1]{%
\vspace{2\baselineskip}
\noindent\textbf{#1}
}

\notebooktitle{Model Evaluation Metrics}
\notebooksubtitle{Classification and Regression Quick Reference}
\notebooktagline{Choose the metric that matches the cost of mistakes.}
\notebookauthor{by Bhuvnesh Verma}
\notebooksource{github.com/MasterBhuvnesh/notebooks}

\begin{document}
\notebookcover

\section{Evaluation Basics}

\field{Good metrics are decision tools.}{A model can look strong on one metric and still fail in production. The best metric is the one that reflects the real business cost of being wrong.}

\begin{tabularx}{\textwidth}{@{}lX@{}}
\toprule
\rowcolor{tablehead} \textbf{Question} & \textbf{What to look for} \\
\midrule
Is the target imbalanced? & Accuracy is often misleading. Check precision, recall, F1, or PR-AUC. \\
Are false negatives costly? & Focus on recall and sensitivity. \\
Are false positives costly? & Focus on precision and specificity. \\
Are predictions continuous? & Use MAE, MSE, RMSE, or R\textsuperscript{2}. \\
Do you need ranking quality? & Use ROC curve or PR curve. \\
\bottomrule
\end{tabularx}

\newpage
\section{Classification Metrics}

\subsection{Confusion Matrix}

\field{What is it?}{A confusion matrix counts correct and incorrect predictions by class. It is the clearest way to see where a classifier fails.}

\field{Used For:}{Classification}

\field{When to use:}{Diagnosing error types, checking class imbalance, fixing recall or precision issues, and understanding the model's mistakes.}

\field{Formula:}{For binary classification, the matrix is
\[
\begin{bmatrix}
TP & FP \\
FN & TN
\end{bmatrix}
\]
where each cell is a count of predictions.}

\field{Symbols:}{\(TP\) = true positives, \(FP\) = false positives, \(FN\) = false negatives, \(TN\) = true negatives.}

\field{What Does It Measure?}{It measures the distribution of predictions across the actual classes. It tells you whether the model is confusing one class for another.}

\field{Interpretation:}{
\valuetable{
Large diagonal, small off-diagonal & Good separation between classes. \\
\midrule
Many \(FN\) & The model misses positive cases. \\
\midrule
Many \(FP\) & The model is too eager to label positives. \\
}}

\field{Example:}{Suppose 100 patients are tested for a rare disease. Results are: \(TP = 24\), \(FN = 6\), \(FP = 8\), \(TN = 62\). The confusion matrix is:
\[
\begin{bmatrix}
24 & 8 \\
6 & 62
\end{bmatrix}
\]
This means 24 diseased patients were correctly detected, 6 were missed, 8 healthy patients were wrongly flagged, and 62 healthy patients were correctly rejected.}

\pyheading{Python (Using Library):}

\begin{py}
from sklearn.metrics import confusion_matrix

y_true = [1, 0, 1, 1, 0, 0]
y_pred = [1, 0, 0, 1, 0, 1]

cm = confusion_matrix(y_true, y_pred)
print(cm)
# [[2 1]
#  [1 2]]
\end{py}

\pyheading{Python (Without Library):}

\begin{py}
def confusion_matrix_py(y_true, y_pred):
    tp = fp = fn = tn = 0
    for yt, yp in zip(y_true, y_pred):
        if yt == 1 and yp == 1:
            tp += 1
        elif yt == 0 and yp == 1:
            fp += 1
        elif yt == 1 and yp == 0:
            fn += 1
        else:
            tn += 1
    return [[tp, fp], [fn, tn]]
\end{py}

\pyheading{Quick Notes:}
\begin{itemize}
\item Advantages: shows exact error pattern; useful for class imbalance and threshold tuning.
\item Limitations: not a single score; can be hard to compare across models at a glance.
\item Best use cases: binary classification, medical screening, fraud detection, churn prediction.
\end{itemize}

\subsection{Accuracy}

\field{What is it?}{Accuracy is the share of all predictions that are correct.}

\field{Used For:}{Classification}

\field{When to use:}{Balanced datasets with roughly equal class frequencies and similar error costs.}

\field{Formula:}{
\[
\text{Accuracy} = \frac{TP + TN}{TP + TN + FP + FN}
\]
}

\field{Symbols:}{\(TP\), \(TN\), \(FP\), \(FN\) are the confusion-matrix counts.}

\field{What Does It Measure?}{It measures overall correctness, but it does not tell you whether the model is failing on the important class.}

\field{Interpretation:}{
\valuetable{
Near 1.0 & Very accurate overall. \\
\midrule
Near 0.5 & Barely better than random guessing. \\
\midrule
Lower than expected on imbalanced data & High accuracy may hide poor minority-class performance. \\
}}

\field{Example:}{Using the previous confusion matrix, accuracy is
\[
\frac{24 + 62}{24 + 62 + 8 + 6} = \frac{86}{100} = 0.86
\]
So the model is 86\% accurate overall.}

\pyheading{Python (Using Library):}

\begin{py}
from sklearn.metrics import accuracy_score

y_true = [1, 0, 1, 1, 0, 0]
y_pred = [1, 0, 0, 1, 0, 1]
print(accuracy_score(y_true, y_pred))
# 0.6666666666666666
\end{py}

\pyheading{Python (Without Library):}

\begin{py}
def accuracy_score_py(y_true, y_pred):
    correct = sum(1 for yt, yp in zip(y_true, y_pred) if yt == yp)
    return correct / len(y_true)
\end{py}

\pyheading{Quick Notes:}
\begin{itemize}
\item Advantages: easy to understand and communicate.
\item Limitations: can be misleading on imbalanced data.
\item Best use cases: balanced classification tasks and baseline comparisons.
\end{itemize}

\subsection{Precision}

\field{What is it?}{Precision asks: of the instances predicted as positive, how many were actually positive?}

\field{Used For:}{Classification}

\field{When to use:}{When false positives are costly, such as email spam detection or fraud alerts.}

\field{Formula:}{
\[
\text{Precision} = \frac{TP}{TP + FP}
\]
}

\field{Symbols:}{\(TP\) = true positives, \(FP\) = false positives.}

\field{What Does It Measure?}{It measures the quality of the positive predictions. A high precision means the model is not labelling too many negatives as positives.}

\field{Interpretation:}{
\valuetable{
Closer to 1.0 & More reliable positive predictions. \\
\midrule
Closer to 0.0 & Many false alarms. \\
}}

\field{Example:}{Using the confusion matrix above,
\[
\text{Precision} = \frac{24}{24 + 8} = \frac{24}{32} = 0.75
\]
So 75\% of predicted positive cases were truly positive.}

\pyheading{Python (Using Library):}

\begin{py}
from sklearn.metrics import precision_score

y_true = [1, 0, 1, 1, 0, 0]
y_pred = [1, 0, 0, 1, 0, 1]
print(precision_score(y_true, y_pred))
# 0.6666666666666666
\end{py}

\pyheading{Python (Without Library):}

\begin{py}
def precision_score_py(y_true, y_pred):
    tp = sum(1 for yt, yp in zip(y_true, y_pred) if yt == 1 and yp == 1)
    fp = sum(1 for yt, yp in zip(y_true, y_pred) if yt == 0 and yp == 1)
    return tp / (tp + fp) if (tp + fp) else 0.0
\end{py}

\pyheading{Quick Notes:}
\begin{itemize}
\item Advantages: useful when false positives are expensive.
\item Limitations: ignores false negatives.
\item Best use cases: fraud detection, email filtering, spam detection.
\end{itemize}

\subsection{Recall}

\field{What is it?}{Recall asks: of the actual positive cases, how many did the model find?}

\field{Used For:}{Classification}

\field{When to use:}{When missing positive cases is costly, such as disease detection or churn prevention.}

\field{Formula:}{
\[
\text{Recall} = \frac{TP}{TP + FN}
\]
}

\field{Symbols:}{\(TP\) = true positives, \(FN\) = false negatives.}

\field{What Does It Measure?}{It measures sensitivity to the positive class. A high recall means the model catches most real positives.}

\field{Interpretation:}{
\valuetable{
Closer to 1.0 & Few positives are missed. \\
\midrule
Closer to 0.0 & Many real positives are being missed. \\
}}

\field{Example:}{With the same confusion matrix,
\[
\text{Recall} = \frac{24}{24 + 6} = \frac{24}{30} = 0.80
\]
The model finds 80\% of actual positive cases.}

\pyheading{Python (Using Library):}

\begin{py}
from sklearn.metrics import recall_score

y_true = [1, 0, 1, 1, 0, 0]
y_pred = [1, 0, 0, 1, 0, 1]
print(recall_score(y_true, y_pred))
# 0.8
\end{py}

\pyheading{Python (Without Library):}

\begin{py}
def recall_score_py(y_true, y_pred):
    tp = sum(1 for yt, yp in zip(y_true, y_pred) if yt == 1 and yp == 1)
    fn = sum(1 for yt, yp in zip(y_true, y_pred) if yt == 1 and yp == 0)
    return tp / (tp + fn) if (tp + fn) else 0.0
\end{py}

\pyheading{Quick Notes:}
\begin{itemize}
\item Advantages: captures missed positive cases.
\item Limitations: ignores false positives.
\item Best use cases: disease detection, anomaly detection, retention outreach.
\end{itemize}

\subsection{F1 Score}

\field{What is it?}{The F1 score is the harmonic mean of precision and recall. It balances both in one number.}

\field{Used For:}{Classification}

\field{When to use:}{Class imbalance, where accuracy is not informative and both false alarms and missed positives matter.}

\field{Formula:}{
\[
F_1 = 2 \cdot \frac{\text{Precision} \cdot \text{Recall}}{\text{Precision} + \text{Recall}}
\]
}

\field{Symbols:}{\(\text{Precision}\) = fraction of predicted positives that are correct; \(\text{Recall}\) = fraction of actual positives that are found.}

\field{What Does It Measure?}{It measures the trade-off between precision and recall in a single score. It punishes weak performance on either side.}

\field{Interpretation:}{
\valuetable{
1.0 & Perfect precision and recall. \\
\midrule
0.0 & No positive predictions are correct. \\
\midrule
Higher is better & Better balance between catching positives and avoiding false alarms. \\
}}

\field{Example:}{Using precision \(= 0.75\) and recall \(= 0.80\),
\[
F_1 = 2 \cdot \frac{0.75 \cdot 0.80}{0.75 + 0.80} = 2 \cdot \frac{0.60}{1.55} \approx 0.774
\]
The model's F1 score is about 0.77.}

\pyheading{Python (Using Library):}

\begin{py}
from sklearn.metrics import f1_score

y_true = [1, 0, 1, 1, 0, 0]
y_pred = [1, 0, 0, 1, 0, 1]
print(f1_score(y_true, y_pred))
# 0.6666666666666666
\end{py}

\pyheading{Python (Without Library):}

\begin{py}
def f1_score_py(y_true, y_pred):
    tp = sum(1 for yt, yp in zip(y_true, y_pred) if yt == 1 and yp == 1)
    fp = sum(1 for yt, yp in zip(y_true, y_pred) if yt == 0 and yp == 1)
    fn = sum(1 for yt, yp in zip(y_true, y_pred) if yt == 1 and yp == 0)
    precision = tp / (tp + fp) if (tp + fp) else 0.0
    recall = tp / (tp + fn) if (tp + fn) else 0.0
    return 2 * precision * recall / (precision + recall) if (precision + recall) else 0.0
\end{py}

\pyheading{Quick Notes:}
\begin{itemize}
\item Advantages: one-number summary of precision-recall trade-off.
\item Limitations: depends on the chosen positive class and threshold.
\item Best use cases: imbalanced classification, anomaly detection, rare-event prediction.
\end{itemize}

\subsection{Specificity}

\field{What is it?}{Specificity measures how well the model identifies negative cases. It is also called the true negative rate.}

\field{Used For:}{Classification}

\field{When to use:}{When avoiding false alarms matters as much as capturing positives.}

\field{Formula:}{
\[
\text{Specificity} = \frac{TN}{TN + FP}
\]
}

\field{Symbols:}{\(TN\) = true negatives, \(FP\) = false positives.}

\field{What Does It Measure?}{It tells you how often the model correctly says "no" for the negative class.}

\field{Interpretation:}{
\valuetable{
Closer to 1.0 & Very good at ruling out negatives. \\
\midrule
Closer to 0.0 & Many healthy cases are incorrectly flagged. \\
}}

\field{Example:}{With the matrix above,
\[
\text{Specificity} = \frac{62}{62 + 8} = \frac{62}{70} = 0.886
\]
So the model correctly rejects 88.6\% of healthy cases.}

\pyheading{Python (Using Library):}

\begin{py}
from sklearn.metrics import confusion_matrix

y_true = [1, 0, 1, 1, 0, 0]
y_pred = [1, 0, 0, 1, 0, 1]
cm = confusion_matrix(y_true, y_pred)
tn, fp, fn, tp = cm.ravel()
print((tn / (tn + fp)))
# 0.5
\end{py}

\pyheading{Python (Without Library):}

\begin{py}
def specificity_score_py(y_true, y_pred):
    tn = sum(1 for yt, yp in zip(y_true, y_pred) if yt == 0 and yp == 0)
    fp = sum(1 for yt, yp in zip(y_true, y_pred) if yt == 0 and yp == 1)
    return tn / (tn + fp) if (tn + fp) else 0.0
\end{py}

\pyheading{Quick Notes:}
\begin{itemize}
\item Advantages: useful when false positives are expensive.
\item Limitations: less informative alone; it should be read alongside recall or sensitivity.
\item Best use cases: medical screening, quality control, safety systems.
\end{itemize}

\subsection{ROC Curve}

\field{What is it?}{The ROC curve plots true positive rate against false positive rate at different decision thresholds.}

\field{Used For:}{Classification}

\field{When to use:}{Comparing classifiers across thresholds, especially when positive class is rare but ranking matters.}

\field{Formula:}{
\[
TPR = \frac{TP}{TP + FN}, \qquad FPR = \frac{FP}{FP + TN}
\]
The ROC curve is the plot of \(TPR\) against \(FPR\) over different thresholds.}

\field{Symbols:}{\(TPR\) = true positive rate, also recall; \(FPR\) = false positive rate, the share of negatives wrongly predicted as positive.}

\field{What Does It Measure?}{It measures the model's ability to rank positive cases above negative cases as the threshold changes.}

\field{Interpretation:}{
\valuetable{
Curve near upper-left corner & Strong discrimination. \\
\midrule
Diagonal line & Random guessing. \\
\midrule
Lower curve & Worse ranking ability. \\
}}

\field{Example:}{If a model is very good at ranking positives ahead of negatives, its ROC curve will rise sharply near the top-left, giving high true-positive rate while keeping false-positive rate low.}

\pyheading{Python (Using Library):}

\begin{py}
from sklearn.metrics import roc_curve
import numpy as np

proba = np.array([0.95, 0.80, 0.40, 0.30, 0.70, 0.20])
y_true = np.array([1, 1, 0, 0, 1, 0])

fpr, tpr, thresholds = roc_curve(y_true, proba)
print(fpr)
print(tpr)
\end{py}

\pyheading{Python (Without Library):}

\begin{py}
def roc_curve_py(y_true, scores):
    thresholds = sorted(set(scores), reverse=True)
    points = []
    for t in thresholds:
        pred = [1 if s >= t else 0 for s in scores]
        tp = sum(1 for yt, yp in zip(y_true, pred) if yt == 1 and yp == 1)
        fp = sum(1 for yt, yp in zip(y_true, pred) if yt == 0 and yp == 1)
        fn = sum(1 for yt, yp in zip(y_true, pred) if yt == 1 and yp == 0)
        tn = sum(1 for yt, yp in zip(y_true, pred) if yt == 0 and yp == 0)
        tpr = tp / (tp + fn) if (tp + fn) else 0.0
        fpr = fp / (fp + tn) if (fp + tn) else 0.0
        points.append((fpr, tpr, t))
    return points
\end{py}

\pyheading{Quick Notes:}
\begin{itemize}
\item Advantages: threshold-independent view of ranking quality.
\item Limitations: less informative when class imbalance is extreme and the operational threshold is fixed.
\item Best use cases: comparing candidate models, ranking systems, probability calibration checks.
\end{itemize}

\subsection{AUC}

\field{What is it?}{AUC is the area under the ROC curve. It summarizes the ROC curve into a single number.}

\field{Used For:}{Classification}

\field{When to use:}{Comparing the overall ranking power of models when the decision threshold is not fixed.}

\field{Formula:}{
\[
\text{AUC} = \int_0^1 TPR(FPR)\,d(FPR)
\]
}

\field{Symbols:}{\(TPR\) = true positive rate, \(FPR\) = false positive rate.}

\field{What Does It Measure?}{It measures how often the model ranks a random positive example above a random negative example.}

\field{Interpretation:}{
\valuetable{
1.0 & Perfect ranking. \\
\midrule
0.5 & Random guessing. \\
\midrule
Below 0.5 & Worse than random. \\
\midrule
Higher is better & Better model discrimination. \\
}}

\field{Example:}{If the ROC curve is close to the upper-left corner, the AUC may be 0.90 or 0.95. That means the model is often correct when ordering a random positive versus a random negative case.}

\pyheading{Python (Using Library):}

\begin{py}
from sklearn.metrics import roc_auc_score
import numpy as np

proba = np.array([0.95, 0.80, 0.40, 0.30, 0.70, 0.20])
y_true = np.array([1, 1, 0, 0, 1, 0])
print(roc_auc_score(y_true, proba))
\end{py}

\pyheading{Python (Without Library):}

\begin{py}
def auc_score_py(y_true, scores):
    pos = [s for yt, s in zip(y_true, scores) if yt == 1]
    neg = [s for yt, s in zip(y_true, scores) if yt == 0]
    total = 0
    for p in pos:
        for n in neg:
            total += 1 if p > n else 0.5 if p == n else 0
    return total / (len(pos) * len(neg)) if pos and neg else 0.0
\end{py}

\pyheading{Quick Notes:}
\begin{itemize}
\item Advantages: single scalar summary of ranking quality; threshold-independent.
\item Limitations: does not show where the operating threshold lies.
\item Best use cases: credit scoring, fraud models, recommendation ranking, binary classification benchmarks.
\end{itemize}

\subsection{Precision-Recall Curve}

\field{What is it?}{The precision-recall curve shows how precision and recall trade off as the decision threshold changes.}

\field{Used For:}{Classification}

\field{When to use:}{Class imbalance or when the positive class is rare and precision matters a lot.}

\field{Formula:}{
\[
\text{Precision}(t) = \frac{TP(t)}{TP(t)+FP(t)},\qquad \text{Recall}(t)=\frac{TP(t)}{TP(t)+FN(t)}
\]
with threshold \(t\) varying across the score range.}

\field{Symbols:}{\(t\) = decision threshold.}

\field{What Does It Measure?}{It measures whether the model can keep precision high while still finding most positive cases.}

\field{Interpretation:}{
\valuetable{
Curve stays high & Good trade-off between precision and recall. \\
\midrule
Sharp drop in precision at high recall & Too many false alarms when recall rises. \\
\midrule
Higher area under the curve & Better performance on rare positives. \\
}}

\field{Example:}{In fraud detection, you may accept lower recall if precision is critical. The PR curve lets you inspect that trade-off directly.}

\pyheading{Python (Using Library):}

\begin{py}
from sklearn.metrics import precision_recall_curve
import numpy as np

proba = np.array([0.95, 0.90, 0.70, 0.30, 0.60, 0.10])
y_true = np.array([1, 1, 0, 0, 1, 0])
precision, recall, thresholds = precision_recall_curve(y_true, proba)
print(precision)
print(recall)
\end{py}

\pyheading{Python (Without Library):}

\begin{py}
def precision_recall_curve_py(y_true, scores):
    thresholds = sorted(set(scores), reverse=True)
    results = []
    for t in thresholds:
        pred = [1 if s >= t else 0 for s in scores]
        tp = sum(1 for yt, yp in zip(y_true, pred) if yt == 1 and yp == 1)
        fp = sum(1 for yt, yp in zip(y_true, pred) if yt == 0 and yp == 1)
        fn = sum(1 for yt, yp in zip(y_true, pred) if yt == 1 and yp == 0)
        precision = tp / (tp + fp) if (tp + fp) else 0.0
        recall = tp / (tp + fn) if (tp + fn) else 0.0
        results.append((t, precision, recall))
    return results
\end{py}

\pyheading{Quick Notes:}
\begin{itemize}
\item Advantages: very informative for rare positives and class imbalance.
\item Limitations: less standard than ROC in some business settings.
\item Best use cases: fraud, default prediction, disease detection, anomaly detection.
\end{itemize}

\subsection{Log Loss}

\field{What is it?}{Log loss measures how close predicted probabilities are to the true class labels. It is a probabilistic loss function.}

\field{Used For:}{Classification}

\field{When to use:}{Comparing probabilistic predictions, especially in binary or multiclass classification.}

\field{Formula:}{
\[
\text{LogLoss} = -\frac{1}{n}\sum_{i=1}^{n}\left[y_i\log(p_i)+(1-y_i)\log(1-p_i)\right]
\]
}

\field{Symbols:}{\(n\) = number of examples; \(y_i\in\{0,1\}\) = true label; \(p_i\) = predicted probability of class 1.}

\field{What Does It Measure?}{It penalizes confident wrong predictions more heavily than modest mistakes. The model is rewarded for probabilities close to the truth.}

\field{Interpretation:}{
\valuetable{
Lower is better & Better probability calibration and predictions. \\
\midrule
Very high values & The model is confident but wrong. \\
}}

\field{Example:}{Suppose the true label is 1 and the model predicts \(p=0.9\). Then the loss is
\[
-\log(0.9) \approx 0.105
\]
If the model predicts \(p=0.2\), the loss is
\[
-\log(0.2) \approx 1.609
\]
The second prediction is much worse, even though it is still a class prediction.}

\pyheading{Python (Using Library):}

\begin{py}
from sklearn.metrics import log_loss

proba = [0.9, 0.2, 0.8, 0.3]
y_true = [1, 0, 1, 0]
print(log_loss(y_true, proba))
\end{py}

\pyheading{Python (Without Library):}

\begin{py}
import math

def log_loss_py(y_true, probs):
    total = 0.0
    for y, p in zip(y_true, probs):
        p = min(max(p, 1e-15), 1 - 1e-15)
        total += -(y * math.log(p) + (1 - y) * math.log(1 - p))
    return total / len(y_true)
\end{py}

\pyheading{Quick Notes:}
\begin{itemize}
\item Advantages: evaluates probability quality, not just hard labels; sensitive to confident mistakes.
\item Limitations: not directly interpretable as accuracy or business impact; harder to explain to non-technical stakeholders.
\item Best use cases: model comparison, probability predictions, ranking and calibration tasks.
\end{itemize}

\newpage
\section{Regression Metrics}

\subsection{MAE (Mean Absolute Error)}

\field{What is it?}{MAE measures the average absolute difference between predicted and actual values.}

\field{Used For:}{Regression}

\field{When to use:}{When you want an easy-to-explain average error in the same units as the target.}

\field{Formula:}{
\[
\text{MAE} = \frac{1}{n}\sum_{i=1}^{n}|y_i - \hat{y}_i|
\]
}

\field{Symbols:}{\(n\) = number of samples; \(y_i\) = actual value; \(\hat{y}_i\) = predicted value.}

\field{What Does It Measure?}{It measures how far predictions are from the truth, on average, without squaring large errors.}

\field{Interpretation:}{
\valuetable{
Lower is better & Predictions are closer to the actual values. \\
}}

\field{Example:}{Actual values: \([10, 12, 14]\); predictions: \([12, 11, 13]\).
\[
|10-12| + |12-11| + |14-13| = 2 + 1 + 1 = 4
\]
\[
\text{MAE} = \frac{4}{3} \approx 1.33
\]
The model misses by about 1.33 units on average.}

\pyheading{Python (Using Library):}

\begin{py}
from sklearn.metrics import mean_absolute_error

y_true = [10, 12, 14]
y_pred = [12, 11, 13]
print(mean_absolute_error(y_true, y_pred))
# 1.3333333333333333
\end{py}

\pyheading{Python (Without Library):}

\begin{py}
def mae_py(y_true, y_pred):
    return sum(abs(y - p) for y, p in zip(y_true, y_pred)) / len(y_true)
\end{py}

\pyheading{Quick Notes:}
\begin{itemize}
\item Advantages: easy to interpret; robust to outliers compared with MSE.
\item Limitations: does not penalize large errors as strongly as squared-error metrics.
\item Best use cases: sales forecasting, household energy prediction, average error reporting.
\end{itemize}

\subsection{MSE (Mean Squared Error)}

\field{What is it?}{MSE is the average of squared prediction errors.}

\field{Used For:}{Regression}

\field{When to use:}{When larger errors should be punished more strongly than smaller ones.}

\field{Formula:}{
\[
\text{MSE} = \frac{1}{n}\sum_{i=1}^{n}(y_i - \hat{y}_i)^2
\]
}

\field{Symbols:}{\(y_i\) = actual value, \(\hat{y}_i\) = predicted value.}

\field{What Does It Measure?}{It measures the average squared distance between predictions and actual values. Squaring emphasizes large mistakes.}

\field{Interpretation:}{
\valuetable{
Lower is better & Predictions are closer overall. \\
\midrule
Large values & Big errors are dominating the score. \\
}}

\field{Example:}{Using the same values,
\[
(10-12)^2 + (12-11)^2 + (14-13)^2 = 4 + 1 + 1 = 6
\]
\[
\text{MSE} = \frac{6}{3} = 2
\]
The model has squared error of 2.}

\pyheading{Python (Using Library):}

\begin{py}
from sklearn.metrics import mean_squared_error

y_true = [10, 12, 14]
y_pred = [12, 11, 13]
print(mean_squared_error(y_true, y_pred))
# 2.0
\end{py}

\pyheading{Python (Without Library):}

\begin{py}
def mse_py(y_true, y_pred):
    return sum((y - p) ** 2 for y, p in zip(y_true, y_pred)) / len(y_true)
\end{py}

\pyheading{Quick Notes:}
\begin{itemize}
\item Advantages: penalizes large errors strongly; widely used in optimization.
\item Limitations: units are squared, so harder to interpret directly.
\item Best use cases: optimization objectives, regression benchmarking, loss functions.
\end{itemize}

\subsection{RMSE (Root Mean Squared Error)}

\field{What is it?}{RMSE is the square root of MSE, bringing the error back to the original unit scale.}

\field{Used For:}{Regression}

\field{When to use:}{When you want a metric in the same units as the target while still penalizing large errors.}

\field{Formula:}{
\[
\text{RMSE} = \sqrt{\frac{1}{n}\sum_{i=1}^{n}(y_i - \hat{y}_i)^2}
\]
}

\field{Symbols:}{\(n\) = sample count; \(y_i\) = actual value; \(\hat{y}_i\) = predicted value.}

\field{What Does It Measure?}{It measures the typical size of the residuals in the original units. It gives more weight to large mistakes than MAE.}

\field{Interpretation:}{
\valuetable{
Lower is better & Error magnitude is smaller. \\
\midrule
Large differences between RMSE and MAE & Large errors are present. \\
}}

\field{Example:}{Since MSE is 2,
\[
\text{RMSE} = \sqrt{2} \approx 1.41
\]
The typical error is about 1.41 units.}

\pyheading{Python (Using Library):}

\begin{py}
from sklearn.metrics import mean_squared_error
import math

y_true = [10, 12, 14]
y_pred = [12, 11, 13]
rmse = math.sqrt(mean_squared_error(y_true, y_pred))
print(rmse)
# 1.4142135623730951
\end{py}

\pyheading{Python (Without Library):}

\begin{py}
import math

def rmse_py(y_true, y_pred):
    mse = sum((y - p) ** 2 for y, p in zip(y_true, y_pred)) / len(y_true)
    return math.sqrt(mse)
\end{py}

\pyheading{Quick Notes:}
\begin{itemize}
\item Advantages: easy to interpret; widely used in forecasting.
\item Limitations: still sensitive to outliers.
\item Best use cases: demand forecasting, house-price prediction, risk models.
\end{itemize}

\subsection{R\textsuperscript{2} Score}

\field{What is it?}{R\textsuperscript{2} measures how much better the model is than predicting the mean of the target.}

\field{Used For:}{Regression}

\field{When to use:}{Comparing regression models and checking how much variance is explained by the model.}

\field{Formula:}{
\[
R^2 = 1 - \frac{\sum_{i=1}^{n}(y_i - \hat{y}_i)^2}{\sum_{i=1}^{n}(y_i - \bar{y})^2}
\]
}

\field{Symbols:}{\(\bar{y}\) = mean of actual target values; \(y_i\) = actual value; \(\hat{y}_i\) = predicted value.}

\field{What Does It Measure?}{It compares the model's residual error to the baseline error of always predicting the mean. Higher values mean the model explains more variance.}

\field{Interpretation:}{
\valuetable{
1.0 & Perfect predictions. \\
\midrule
0.0 & Model performs like using the mean baseline. \\
\midrule
Negative values & Worse than the mean baseline. \\
\midrule
Higher is better & Better explanatory power. \\
}}

\field{Example:}{Suppose actual values are \([10, 12, 14]\) and predictions are \([12, 11, 13]\). The mean is 12. The numerator is 6 and the denominator is 4. So:
\[
R^2 = 1 - \frac{6}{4} = -0.5
\]
This is worse than predicting the mean.}

\pyheading{Python (Using Library):}

\begin{py}
from sklearn.metrics import r2_score

y_true = [10, 12, 14]
y_pred = [12, 11, 13]
print(r2_score(y_true, y_pred))
# -0.5
\end{py}

\pyheading{Python (Without Library):}

\begin{py}
def r2_score_py(y_true, y_pred):
    y_mean = sum(y_true) / len(y_true)
    ss_res = sum((y - p) ** 2 for y, p in zip(y_true, y_pred))
    ss_tot = sum((y - y_mean) ** 2 for y in y_true)
    return 1 - (ss_res / ss_tot) if ss_tot else 0.0
\end{py}

\pyheading{Quick Notes:}
\begin{itemize}
\item Advantages: interpretable and common in regression reporting.
\item Limitations: can increase when adding features even if they are not useful; does not penalize complexity.
\item Best use cases: model comparison, variance explained, regression summaries.
\end{itemize}

\subsection{Adjusted R\textsuperscript{2}}

\field{What is it?}{Adjusted R\textsuperscript{2} modifies R\textsuperscript{2} to account for the number of predictors in the model.}

\field{Used For:}{Regression}

\field{When to use:}{Comparing models with different numbers of variables, especially in feature selection.}

\field{Formula:}{
\[
\text{Adjusted } R^2 = 1 - \left(\frac{1-R^2}{n-k-1}\right)(n-1)
\]
}

\field{Symbols:}{\(n\) = sample size; \(k\) = number of predictors; \(R^2\) = coefficient of determination.}

\field{What Does It Measure?}{It tells you whether added predictors improve the model enough to justify their complexity.}

\field{Interpretation:}{
\valuetable{
Higher is better & Better balance between fit and complexity. \\
\midrule
Can be lower than R\textsuperscript{2} & Extra variables are not adding value. \\
}}

\field{Example:}{If a model has \(R^2 = 0.80\), \(n = 100\), and \(k = 5\), then
\[
\text{Adjusted } R^2 = 1 - \left(\frac{1 - 0.80}{100 - 5 - 1}\right)(99)
\approx 0.79
\]
The adjusted value is slightly lower than the raw R\textsuperscript{2}, reflecting model complexity.}

\pyheading{Python (Using Library):}

\begin{py}
from sklearn.linear_model import LinearRegression
from sklearn.metrics import r2_score
import numpy as np

X = np.array([[1], [2], [3], [4], [5]])
y = np.array([2, 4, 6, 8, 10])
model = LinearRegression().fit(X, y)
print(r2_score(y, model.predict(X)))
\end{py}

\pyheading{Python (Without Library):}

\begin{py}
def adjusted_r2_score_py(y_true, y_pred, k):
    n = len(y_true)
    r2 = 1 - sum((y - p) ** 2 for y, p in zip(y_true, y_pred)) / sum((y - sum(y_true) / n) ** 2 for y in y_true)
    return 1 - ((1 - r2) * (n - 1) / (n - k - 1)) if n - k - 1 > 0 else r2
\end{py}

\pyheading{Quick Notes:}
\begin{itemize}
\item Advantages: better than raw R\textsuperscript{2} for comparing models with different feature counts.
\item Limitations: still not a universal metric; choose based on the task and data.
\item Best use cases: feature selection, multi-variable regression, model comparison.
\end{itemize}

\subsection{MAPE (Mean Absolute Percentage Error)}

\field{What is it?}{MAPE measures prediction error as a percentage of the actual value.}

\field{Used For:}{Regression}

\field{When to use:}{When you want a relative error metric that is easy to explain in percentage terms.}

\field{Formula:}{
\[
\text{MAPE} = \frac{100}{n}\sum_{i=1}^{n}\left|\frac{y_i - \hat{y}_i}{y_i}\right|
\]
}

\field{Symbols:}{\(y_i\) = actual value, \(\hat{y}_i\) = predicted value, \(n\) = sample count.}

\field{What Does It Measure?}{It measures the average percentage gap between prediction and reality. It is useful for communicating forecast error in familiar percentage terms.}

\field{Interpretation:}{
\valuetable{
Lower is better & Lower percentage error. \\
\midrule
Very high values & Forecasts are often far from actual values. \\
}}

\field{Example:}{Actual values: \([100, 200, 300]\), predictions: \([90, 210, 330]\).
\[
\left|\frac{100-90}{100}\right|=0.10,\quad
\left|\frac{200-210}{200}\right|=0.05,\quad
\left|\frac{300-330}{300}\right|=0.10
\]
\[
\text{MAPE} = \frac{100}{3}(0.10 + 0.05 + 0.10) \approx 8.33\%
\]
The average percentage error is about 8.33\%.}

\pyheading{Python (Using Library):}

\begin{py}
from sklearn.metrics import mean_absolute_percentage_error

y_true = [100, 200, 300]
y_pred = [90, 210, 330]
print(mean_absolute_percentage_error(y_true, y_pred))
# 0.08333333333333334
\end{py}

\pyheading{Python (Without Library):}

\begin{py}
def mape_py(y_true, y_pred):
    total = 0.0
    for y, p in zip(y_true, y_pred):
        if y == 0:
            continue
        total += abs((y - p) / y)
    return 100 * total / len(y_true)
\end{py}

\pyheading{Quick Notes:}
\begin{itemize}
\item Advantages: easy to communicate as percentages; useful for business forecasting.
\item Limitations: breaks down when actual values are near zero; can be skewed by small denominators.
\item Best use cases: time-series forecasting, demand prediction, business KPIs.
\end{itemize}

\section{Practical Rule of Thumb}

\begin{tabularx}{\textwidth}{@{}lX@{}}
\toprule
\rowcolor{tablehead} \textbf{Problem type} & \textbf{Best starting metric} \\
\midrule
Binary classification, balanced classes & Accuracy \\
Rare positive events, false negatives costly & Recall, F1, PR-AUC \\
False positives costly & Precision, specificity \\
Ranking or probability quality & ROC-AUC, log loss \\
Regression with interpretable errors & MAE, RMSE \\
Regression with relative error reporting & MAPE \\
Comparing model fit while adjusting for complexity & Adjusted R\textsuperscript{2} \\
\bottomrule
\end{tabularx}

\field{Important note:}{No single metric is universally best. Choose the metric that matches the consequence of being wrong, then validate with the business or task objective.}

\end{document}
